% ==========================================================
% NOVEL NDVI-BASED VEGETATION STRESS ANALYSIS
% For: Forest Fire Susceptibility Mapping — India (2000–2025)
% Extends: Biswas et al. (2025), Environ. Sci. Pollut. Res., 32:4856–4878
%
% Author : [Your Name], Senior Research Fellow (SRF)
% Dept.  : Applied Mathematics & Scientific Computing
% Institute: IIT Roorkee
% Supervisor: Prof. Millie Pant
% Funding : Ministry of Education, Government of India
%
% HOW TO COMPILE (TeXstudio):
%   Step 1: F6  (pdflatex)
%   Step 2: F8  (bibtex)
%   Step 3: F6  (pdflatex)
%   Step 4: F6  (pdflatex)
%   OR simply press F5 (Build & View)
%
% COMPILER: pdflatex
% ==========================================================

\documentclass[12pt, a4paper]{article}

% ── Core packages ──────────────────────────────────────────
\usepackage[utf8]{inputenc}
\usepackage[T1]{fontenc}
\usepackage[english]{babel}
\usepackage{lmodern}

% ── Page geometry ──────────────────────────────────────────
\usepackage[
  a4paper,
  left   = 2.8cm,
  right  = 2.8cm,
  top    = 2.5cm,
  bottom = 2.5cm
]{geometry}

% ── Mathematics ────────────────────────────────────────────
\usepackage{amsmath}
\usepackage{amssymb}
\usepackage{amsthm}
\usepackage{bm}
\usepackage{mathtools}

% ── Graphics & Tables ──────────────────────────────────────
\usepackage{graphicx}
\usepackage{booktabs}
\usepackage{array}
\usepackage{multirow}
\usepackage{longtable}
\usepackage{tabularx}
\usepackage{float}
\usepackage{colortbl}          % required for \rowcolor in tabular

% ── Colours & Boxes ────────────────────────────────────────
\usepackage{xcolor}
\usepackage[most]{tcolorbox}
\tcbuselibrary{skins, breakable, theorems}

\definecolor{iitblue}{RGB}{26,82,118}
\definecolor{midblue}{RGB}{46,117,182}
\definecolor{lightblue}{RGB}{214,234,248}
\definecolor{darkgreen}{RGB}{20,90,50}
\definecolor{midgreen}{RGB}{30,132,73}
\definecolor{lightgreen}{RGB}{213,245,227}
\definecolor{darkred}{RGB}{146,43,33}
\definecolor{lightred}{RGB}{250,219,216}
\definecolor{purple}{RGB}{108,52,131}
\definecolor{lightpurple}{RGB}{232,218,239}
\definecolor{gold}{RGB}{183,149,11}
\definecolor{lightgold}{RGB}{254,249,231}
\definecolor{gray1}{RGB}{242,242,242}
\definecolor{gray2}{RGB}{217,217,217}

% ── Custom tcolorbox environments ──────────────────────────
\newtcolorbox{novelbox}[1]{
  enhanced, breakable,
  colback   = lightgreen,
  colframe  = darkgreen,
  fonttitle = \bfseries\color{white},
  title     = {\textbf{$\bigstar$\; Novel Contribution:\; #1}},
  attach boxed title to top left = {yshift=-2mm, xshift=4mm},
  boxed title style = {colback=darkgreen, colframe=darkgreen},
  top = 6pt, bottom = 6pt
}

\newtcolorbox{mathbox}[1]{
  enhanced,
  colback   = lightblue,
  colframe  = midblue,
  fonttitle = \bfseries\color{white},
  title     = {#1},
  attach boxed title to top left = {yshift=-2mm, xshift=4mm},
  boxed title style = {colback=iitblue, colframe=iitblue},
  top = 8pt, bottom = 8pt
}

\newtcolorbox{simplebox}[1]{
  enhanced,
  colback   = lightgold,
  colframe  = gold,
  fonttitle = \bfseries\color{gold},
  title     = {Simple Explanation: #1},
  top = 6pt, bottom = 6pt
}

\newtcolorbox{gapbox}{
  enhanced, breakable,
  colback   = lightred,
  colframe  = darkred,
  fonttitle = \bfseries\color{white},
  title     = {Gap in Biswas et al.\ (2025) Addressed by This Study},
  top = 6pt, bottom = 6pt
}

% ── Typography ─────────────────────────────────────────────
\usepackage{setspace}
\onehalfspacing
\usepackage{parskip}
\setlength{\parindent}{0pt}
\setlength{\parskip}{6pt}

% ── References ─────────────────────────────────────────────
\usepackage[round, authoryear, sort]{natbib}

% ── Hyperlinks ─────────────────────────────────────────────
\usepackage[
  colorlinks = true,
  linkcolor  = iitblue,
  citecolor  = darkgreen,
  urlcolor   = midblue
]{hyperref}

% ── Header / Footer ────────────────────────────────────────
\usepackage{fancyhdr}
\pagestyle{fancy}
\fancyhf{}
\rhead{\small\color{midblue} Novel NDVI Methodology $|$ IIT Roorkee}
\lhead{\small\color{midblue} Forest Fire Susceptibility Mapping --- India}
\cfoot{\small\thepage}
\renewcommand{\headrulewidth}{0.5pt}

% ── Theorem environments ───────────────────────────────────
\newtheorem{definition}{Definition}[section]
\newtheorem{proposition}{Proposition}[section]
\newtheorem{remark}{Remark}[section]

% ── Algorithm ──────────────────────────────────────────────
\usepackage{algorithm}
\usepackage{algpseudocode}

% ── Enumitem ───────────────────────────────────────────────
\usepackage{enumitem}

% ── Caption ────────────────────────────────────────────────
\usepackage{caption}
\captionsetup{font=small, labelfont=bf}

% ==========================================================
\begin{document}
% ==========================================================

% ── Title Page ─────────────────────────────────────────────
\begin{titlepage}
\centering
\vspace*{1.5cm}

{\color{iitblue}\rule{\linewidth}{1.5pt}}
\vspace{0.4cm}

{\LARGE\bfseries\color{iitblue}
Novel NDVI-Based Vegetation Stress Analysis\\[0.5em]
for Forest Fire Susceptibility Mapping in India}

\vspace{0.3cm}
{\color{gold}\rule{\linewidth}{1pt}}
\vspace{0.4cm}

{\large\itshape\color{midblue}
Mathematical Formulation $\cdot$ Physical Justification $\cdot$
Novel Contributions}

\vspace{1.2cm}

\begin{tcolorbox}[enhanced, colback=lightblue, colframe=iitblue,
  width=0.88\textwidth, arc=4pt]
\centering
\textbf{Submitted to:} Prof.\ Millie Pant\\[4pt]
Department of Applied Mathematics \& Scientific Computing\\
Indian Institute of Technology Roorkee\\[8pt]
\textbf{Data:} MOD13A3 v061 $|$ 1\,km monthly NDVI $|$
India $|$ 2000--2025 (309 files)\\[4pt]
\textbf{Extends:} \citealt{biswas2025} ---
\textit{Environ.\ Sci.\ Pollut.\ Res.}, 32:4856--4878\\[4pt]
\textbf{Funding:} Ministry of Education, Government of India
\end{tcolorbox}

\vspace{1cm}
{\large\today}
\vspace*{1cm}

{\color{iitblue}\rule{\linewidth}{1.5pt}}
\end{titlepage}

% ── Table of Contents ──────────────────────────────────────
\tableofcontents
\clearpage

% ==========================================================
\section{Background and Motivation}
\label{sec:background}
% ==========================================================

\begin{gapbox}
In \citet{biswas2025}, NDVI is the \textbf{single most important variable}
for Indian forest fire occurrence prediction:
permutation importance $= 22.3\%$, percentage contribution $= 28.4\%$
(highest among all 15 variables, Table~3 therein).
Yet Biswas et al.\ applied NDVI as a single time-averaged composite map
from MOD13C2 v006, with \textbf{no quality screening, no temporal
decomposition, no anomaly extraction, and no pre-fire stress
quantification.}
\end{gapbox}

The present study extends \citet{biswas2025} by applying five
mathematically rigorous transformations to the 25-year monthly NDVI
time series (MOD13A3 v061, 1\,km, 309 monthly files, 2000--2025) to
derive \textbf{seven novel NDVI-based features} (F1--F7) that are
directly connected to the physics of fire ignition.
Table~\ref{tab:features_overview} summarises all features and their
novelty status.

\begin{table}[H]
\centering
\caption{NDVI-derived feature set: this study vs \citet{biswas2025}}
\label{tab:features_overview}
\small
\begin{tabular}{@{}clllc@{}}
\toprule
\textbf{ID} & \textbf{Feature} & \textbf{Core Equation} &
\textbf{Physical Meaning} & \textbf{Novel?} \\
\midrule
F1 & QA-filtered NDVI & Eq.~\eqref{eq:ndvi_scale} &
    Cloud-screened vegetation density & Partial \\
F2 & Climatological mean $\bar{\mu}^{(m)}$ & Eq.~\eqref{eq:clim_mean} &
    Seasonal baseline per month & No \\
\rowcolor{lightgreen}
F3 & Monthly anomaly $\delta$ & Eq.~\eqref{eq:anomaly} &
    Inter-annual vegetation stress & \textbf{Yes} \\
\rowcolor{lightgreen}
F4 & STL trend $T(t)$ & Eq.~\eqref{eq:stl} &
    25-year vegetation change direction & \textbf{Yes} \\
\rowcolor{lightgreen}
F5 & STL residual $R(t)$ & Eq.~\eqref{eq:stl} &
    Sudden stress: drought, fire-scar & \textbf{Yes} \\
\rowcolor{lightgreen}
F6 & Mann--Kendall $\tau$ & Eq.~\eqref{eq:mk_tau} &
    Statistical significance of trend & \textbf{Yes} \\
\rowcolor{lightgreen}
F7 & CVSI$(k^{*})$ & Eq.~\eqref{eq:cvsi} &
    Cumulative pre-fire moisture deficit & \textbf{Yes} \\
\midrule
\multicolumn{3}{l}{\textit{Biswas et al.\ (2025): 1 feature}} &
\multicolumn{2}{r}{\textit{This study: 7 features (5 novel)}} \\
\bottomrule
\end{tabular}
\end{table}

% ==========================================================
\section{NDVI: Definition and Scale Correction}
\label{sec:ndvi_def}
% ==========================================================

\begin{simplebox}{What is NDVI?}
NDVI measures how green and healthy vegetation is.
A healthy forest absorbs red light (for photosynthesis) and reflects
near-infrared light.
Dry, stressed, or burned vegetation does the opposite.
High NDVI ($>0.6$) = dense healthy forest.
Low NDVI ($<0.2$) = sparse, stressed, or burned vegetation.
\end{simplebox}

\begin{definition}[Normalized Difference Vegetation Index]
\label{def:ndvi}
\begin{equation}
  \mathrm{NDVI} =
  \frac{\rho_{\mathrm{NIR}} - \rho_{\mathrm{Red}}}
       {\rho_{\mathrm{NIR}} + \rho_{\mathrm{Red}}}
  \label{eq:ndvi_def}
\end{equation}
where $\rho_{\mathrm{NIR}}$ and $\rho_{\mathrm{Red}}$ denote surface
reflectance in the Near-Infrared and Red spectral bands, respectively.
The index ranges over $[-0.2,\,1.0]$.
\end{definition}

MODIS MOD13A3 stores NDVI as signed 16-bit integers.
The corrected floating-point NDVI for spatial pixel $(i,j)$ at
monthly time step $t$ is:

\begin{mathbox}{Equation F1 — Scale Correction and Fill-Value Masking}
\begin{equation}
  \mathrm{NDVI}_{ij}^{(t)} =
  \begin{cases}
    \mathrm{raw}_{ij}^{(t)} \times 10^{-4} &
      \text{if } -2000 \le \mathrm{raw}_{ij}^{(t)} \le 10000 \\[4pt]
    \mathrm{NaN} & \text{otherwise}
  \end{cases}
  \label{eq:ndvi_scale}
\end{equation}
\textit{Fill value: $-3000$ (replace with NaN). Valid NDVI range after
scaling: $[-0.2,\,1.0]$.}
\end{mathbox}

\subsection*{Quality Assurance (QA) Filtering}

The \texttt{pixel\_reliability} Science Data Set (SDS) from MOD13A3
provides a pixel-quality flag $q_{ij}^{(t)}$:

\begin{equation}
  \mathrm{NDVI}_{ij}^{(t),\,\mathrm{QA}} =
  \begin{cases}
    \mathrm{NDVI}_{ij}^{(t)} & \text{if } q_{ij}^{(t)} \in \{0,\,1\} \\[4pt]
    \mathrm{NaN}             & \text{otherwise}
  \end{cases}
  \label{eq:qa_filter}
\end{equation}

where $q = 0$ (Good data), $q = 1$ (Marginal data, usable),
$q = 2$ (Snow/Ice --- exclude), $q = 3$ (Cloudy --- exclude),
$q = -1$ (Fill/No-data).
\citet{biswas2025} applied \textbf{no} such quality screening.

% ==========================================================
\section{Novel Contribution 1: STL Temporal Decomposition}
\label{sec:stl}
% ==========================================================

\begin{simplebox}{The 3-Component NDVI Video}
Think of your 309 monthly NDVI images as a 25-year video of India's
forests. Each pixel has a time series with 309 values.
STL separates this into three stories:
(1) Is the forest getting greener or browner over 25 years? [Trend]
(2) What is the normal monsoon cycle? [Seasonal]
(3) What unexpected events happened? [Residual — drought, fire-scars]
\end{simplebox}

\subsection{STL Decomposition Model}

The STL (Seasonal and Trend decomposition using LOESS) algorithm
\citep{cleveland1990stl} decomposes the NDVI time series at each pixel
$(i,j)$ additively:

\begin{mathbox}{Equation F4/F5 — STL Decomposition}
\begin{equation}
  \mathrm{NDVI}_{ij}^{(t)} =
  \underbrace{T_{ij}^{(t)}}_{\text{Trend}} +
  \underbrace{S_{ij}^{(t)}}_{\text{Seasonal}} +
  \underbrace{R_{ij}^{(t)}}_{\text{Residual}}
  \label{eq:stl}
\end{equation}
Applied independently to every spatial pixel.
Period $p = 12$ months. Robust fitting: yes (downweights outliers).
Time index $t = 1, 2, \ldots, T$ where $T = 309$ monthly steps.
\end{mathbox}

\subsection{Component Definitions and Physical Interpretations}

\paragraph{Trend Component $T_{ij}^{(t)}$ (Feature F4).}
The trend is estimated via a LOESS (Locally Estimated Scatterplot
Smoothing) smoother.
For a time point $t_0$, the LOESS estimate is obtained by solving the
weighted local polynomial regression:

\begin{equation}
  \hat{T}(t_0) =
  \underset{\beta_0,\,\beta_1}{\arg\min}
  \sum_{t=1}^{T}
  \mathcal{K}_h(t - t_0)
  \left[
    \mathrm{NDVI}_{ij}^{(t)} - \beta_0 - \beta_1(t - t_0)
  \right]^2
  \label{eq:loess}
\end{equation}

where $\mathcal{K}_h$ is the tri-cube kernel weight function:

\begin{equation}
  \mathcal{K}_h(u) =
  \left(1 - \left|\frac{u}{h}\right|^3\right)^3
  \cdot \mathbf{1}\!\left[|u| < h\right]
  \label{eq:tricube}
\end{equation}

and $h > 0$ is the bandwidth controlling smoothness.

\begin{remark}
If $\partial T_{ij} / \partial t < 0$ over the study period,
the pixel's vegetation is in long-term decline --- indicating
forest degradation, reduced canopy moisture, and elevated fire
susceptibility. \citeauthor{biswas2025} cannot detect this
from a single temporal average.
\end{remark}

\paragraph{Seasonal Component $S_{ij}^{(t)}$ (removed from analysis).}
The periodic annual monsoon cycle ($p = 12$ months) is isolated and
\emph{removed} to de-seasonalise the signal.
This prevents normal dry-season dryness from being misclassified as
genuine vegetation stress.

\paragraph{Residual Component $R_{ij}^{(t)}$ (Feature F5).}
After removing trend and seasonal components:

\begin{equation}
  R_{ij}^{(t)} =
  \mathrm{NDVI}_{ij}^{(t)} - T_{ij}^{(t)} - S_{ij}^{(t)}
  \label{eq:residual}
\end{equation}

Large negative residuals ($R_{ij}^{(t)} \ll 0$) indicate sudden,
unexpected vegetation drops --- drought episodes, pest outbreaks,
or post-fire regeneration patterns.
These are the \textbf{pre-fire stress fingerprints} that a
time-averaged NDVI completely conceals.

\subsection{Mann--Kendall Trend Significance Test (Feature F6)}

After computing the pixel-wise trend component $T_{ij}^{(t)}$,
the non-parametric Mann--Kendall test \citep{mann1945,kendall1975}
determines whether the observed trend is statistically significant
or could arise by random chance.

\subsubsection{Test Statistic}

For a pixel time series $\{T^{(1)}, T^{(2)}, \ldots, T^{(n)}\}$
(where $n = T = 309$), define the Mann--Kendall statistic:

\begin{equation}
  S = \sum_{k=1}^{n-1} \sum_{j=k+1}^{n}
  \operatorname{sgn}\!\left(T^{(j)} - T^{(k)}\right)
  \label{eq:mk_S}
\end{equation}

where:
\begin{equation}
  \operatorname{sgn}(x) =
  \begin{cases}
    +1 & \text{if } x > 0 \\
     0 & \text{if } x = 0 \\
    -1 & \text{if } x < 0
  \end{cases}
  \label{eq:sgn}
\end{equation}

\subsubsection{Normalised Kendall's $\tau$}

\begin{mathbox}{Equation F6 — Mann--Kendall $\tau$ Statistic}
\begin{equation}
  \tau = \frac{P - Q}{\dfrac{1}{2}\,n\,(n-1)}
  \label{eq:mk_tau}
\end{equation}
where $P$ = number of concordant pairs (later value $>$ earlier value),
$Q$ = number of discordant pairs (later value $<$ earlier value),
$\tau \in [-1, +1]$.
\end{mathbox}

\subsubsection{Variance and $z$-Score}

For large $n$ without ties:
\begin{equation}
  \operatorname{Var}(S) =
  \frac{n(n-1)(2n+5)}{18}
  \label{eq:mk_var}
\end{equation}

The standardised test statistic is:
\begin{equation}
  z =
  \begin{cases}
    \dfrac{S - 1}{\sqrt{\operatorname{Var}(S)}} & \text{if } S > 0 \\[8pt]
    0                                           & \text{if } S = 0 \\[4pt]
    \dfrac{S + 1}{\sqrt{\operatorname{Var}(S)}} & \text{if } S < 0
  \end{cases}
  \label{eq:mk_z}
\end{equation}

The two-tailed $p$-value is: $p = 2\left[1 - \Phi(|z|)\right]$
where $\Phi$ is the standard normal CDF.

\subsubsection{Decision Rule and Spatial Output Maps}

\begin{equation}
  \text{Pixel classification} =
  \begin{cases}
    \textbf{Significant Browning} & \tau < 0 \text{ and } p < 0.05 \\
    \textbf{Significant Greening} & \tau > 0 \text{ and } p < 0.05 \\
    \textbf{No Reliable Trend}    & p \ge 0.05
  \end{cases}
  \label{eq:mk_decision}
\end{equation}

The code produces two output maps:
\begin{itemize}[leftmargin=2em]
  \item \texttt{browning\_sig}: pixels with statistically
        significant vegetation decline ($\tau < 0$, $p < 0.05$)
        --- these are the \textbf{highest fire-risk zones}.
  \item \texttt{greening\_sig}: pixels with statistically
        significant vegetation recovery ($\tau > 0$, $p < 0.05$).
\end{itemize}

\begin{novelbox}{Novel Contribution 1 --- STL Decomposition and Mann--Kendall Test}
\begin{itemize}[leftmargin=1.5em,itemsep=2pt]
  \item \citet{biswas2025} used a single time-averaged NDVI value
        from MOD13C2 v006, with no temporal decomposition.
  \item This study applies pixel-wise STL to 309 monthly values
        (2000--2025), producing three physically distinct components:
        Trend $T$ (Feature F4), Seasonal $S$ (removed), Residual $R$
        (Feature F5).
  \item Mann--Kendall $\tau$ (Feature F6) confirms statistical
        significance of each browning/greening pixel.
  \item A pixel with $\tau = -0.31$, $p = 0.002$ in Uttarakhand
        (pauri Garhwal, 30.2\textdegree N, 78.8\textdegree E)
        represents a net NDVI loss of $0.075$ over 25 years ---
        invisible to a single-value approach.
  \item \textbf{References:} \citet{cleveland1990stl};
        \citet{mann1945}; \citet{kendall1975}.
\end{itemize}
\end{novelbox}

% ==========================================================
\section{Novel Contribution 2: Monthly NDVI Anomaly}
\label{sec:anomaly}
% ==========================================================

\begin{simplebox}{Why Remove the Seasonal Pattern?}
Indian forests are always greener in the monsoon (July--September)
and browner in pre-monsoon (March--May). If you look at raw NDVI in
May, every forest looks stressed --- but that is \emph{normal}.
The anomaly $\delta$ strips this away and asks:
``Is this May more stressed than the average May over 20 years?''
A negative $\delta$ means genuine abnormal stress --- the real
fire-relevant signal.
\end{simplebox}

\subsection{Climatological Monthly Mean (Feature F2)}

Using the 2001--2020 baseline (consistent with \citealt{biswas2025}):

\begin{mathbox}{Equation F2 — Climatological Mean}
\begin{equation}
  \bar{\mu}_{ij}^{(m)} =
  \frac{1}{\lvert\mathcal{Y}_{m}\rvert}
  \sum_{y \in \mathcal{Y}_{m}}
  \mathrm{NDVI}_{ij}^{(y,m),\,\mathrm{QA}}
  \label{eq:clim_mean}
\end{equation}
where $m \in \{1, 2, \ldots, 12\}$ is the calendar month,
$\mathcal{Y}_{m} = \{2001, 2002, \ldots, 2020\}$,
and only non-\texttt{NaN} values contribute.
This produces \textbf{12 mean maps} --- one per calendar month.
\end{mathbox}

\subsection{Monthly NDVI Anomaly (Feature F3)}

\begin{mathbox}{Equation F3 — NDVI Anomaly $\delta$}
\begin{equation}
  \boxed{
    \delta_{ij}^{(y,m)} =
    \mathrm{NDVI}_{ij}^{(y,m),\,\mathrm{QA}} - \bar{\mu}_{ij}^{(m)}
  }
  \label{eq:anomaly}
\end{equation}
\begin{itemize}
  \item $\delta_{ij}^{(y,m)} < 0$: below-average greenness
        (\textbf{vegetation stress}) --- fire risk \textbf{increases}.
  \item $\delta_{ij}^{(y,m)} > 0$: above-average greenness
        --- fire risk decreases.
  \item By construction: $\sum_{y \in \mathcal{Y}_m}
        \delta_{ij}^{(y,m)} \approx 0$
        (zero mean over baseline).
\end{itemize}
\end{mathbox}

\begin{proposition}[Anomaly as De-seasonalised Signal]
The anomaly $\delta_{ij}^{(y,m)}$ removes the confounding seasonal
cycle and exposes the inter-annual vegetation stress signal that is
causally connected to fire ignition, independent of the time of year.
\end{proposition}

\begin{novelbox}{Novel Contribution 2 --- Monthly NDVI Anomaly}
\begin{itemize}[leftmargin=1.5em,itemsep=2pt]
  \item \citet{biswas2025} cannot distinguish between
        ``normal dry season'' and ``abnormally stressed vegetation.''
  \item The anomaly $\delta_{ij}^{(y,m)}$ isolates genuine
        inter-annual stress signals from the expected monsoon cycle.
  \item \textbf{Concrete example --- Uttarakhand, March 2016
        (El Ni\~{n}o year):}
        Raw NDVI $= 0.38$ (looks the same as any dry March).
        Anomaly $\delta = 0.38 - 0.44 = -0.06$ ---
        the forest is 6 NDVI-units below its 20-year March average.
        This is the genuine drought stress signal.
        \citet{biswas2025} cannot see this.
  \item This anomaly becomes the direct input to the CVSI
        (Section~\ref{sec:cvsi}) and Feature F3 in the ML pipeline.
\end{itemize}
\end{novelbox}

% ==========================================================
\section{Novel Contribution 3: Cumulative Vegetation Stress Index (CVSI)}
\label{sec:cvsi}
% ==========================================================

\begin{simplebox}{The Pre-Fire Memory of Drought}
Consider a forest in Odisha.
In January: NDVI slightly below normal ($\delta = -0.03$).
In February: still below normal ($\delta = -0.05$).
In March: again below ($\delta = -0.04$).
In April: \textbf{fire breaks out.}

\medskip
Biswas et al.\ see only: April NDVI $= 0.35$.
They cannot see that three months of accumulated stress preceded the fire.

The CVSI adds up all negative anomalies in the $k$ months
\emph{before} each fire event --- quantifying the multi-month
moisture deficit that enabled ignition.
\end{simplebox}

\subsection{Mathematical Definition}

Let $\tau_{f}$ denote the calendar month of fire occurrence at pixel
$(i,j)$.
The $k$-month Cumulative Vegetation Stress Index is:

\begin{mathbox}{Equation F7 — Cumulative Vegetation Stress Index (CVSI)}
\begin{equation}
  \boxed{
    \mathrm{CVSI}_{ij}\!\left(\tau_{f},\, k\right) =
    \sum_{t\,=\,\tau_{f} - k}^{\tau_{f} - 1}
    \max\!\left(-\delta_{ij}^{(t)},\; 0\right)
  }
  \label{eq:cvsi}
\end{equation}
The operator $\max(-\delta, 0)$ retains \textbf{only} stress episodes
(negative anomaly months) and discards above-average greenness months.
$k \in \{1, 2, 3, 4, 5, 6\}$ is the accumulation window (months).
\end{mathbox}

\subsection{Expanded Form for $k = 3$}

\begin{equation}
  \mathrm{CVSI}_{ij}(\tau_{f},\,3) =
  \underbrace{\max\!\left(-\delta_{ij}^{(\tau_{f}-3)},\,0\right)}_{\text{3 months before}}
  +\;
  \underbrace{\max\!\left(-\delta_{ij}^{(\tau_{f}-2)},\,0\right)}_{\text{2 months before}}
  +\;
  \underbrace{\max\!\left(-\delta_{ij}^{(\tau_{f}-1)},\,0\right)}_{\text{1 month before}}
  \label{eq:cvsi3}
\end{equation}

Each term is zero if that month had above-average greenness,
and positive if there was stress.
The sum accumulates the total moisture deficit in the pre-fire window.

\subsection{Optimal Lag Selection via Mutual Information}

The optimal accumulation window $k^{*}$ is selected by maximising the
Shannon mutual information between the CVSI and binary fire occurrence
labels $F_{ij} \in \{0, 1\}$:

\begin{equation}
  k^{*} =
  \underset{k\,\in\,\{1,2,3,4,5,6\}}{\arg\max}\;
  I\!\left(\mathrm{CVSI}(\tau_{f},\,k)\;;\;F_{ij}\right)
  \label{eq:mi_lag}
\end{equation}

where the mutual information is:

\begin{equation}
  I(X;\,Y) =
  \sum_{x,\,y}
  p(x,y)\,\log_{2}\!\frac{p(x,y)}{p(x)\,p(y)}
  \label{eq:mutual_info}
\end{equation}

This information-theoretic criterion identifies which lag window most
reduces uncertainty about fire occurrence --- a rigorous alternative
to the fixed-window approaches used in the literature.

\subsection{Physical Justification}

\begin{proposition}[CVSI as Fuel Moisture Deficit Proxy]
Prolonged below-average NDVI ($\delta < 0$) signals:
\begin{enumerate}[leftmargin=2em]
  \item Reduced canopy moisture content (leaf water potential decreasing),
  \item Accelerated fine fuel drying (dead leaf litter moisture $\downarrow$),
  \item Increased bulk density of surface fuels susceptible to ignition.
\end{enumerate}
The CVSI therefore serves as a proxy for the \textbf{pre-fire fuel
moisture deficit} --- a physically established precursor to wildfire
ignition \citep{chuvieco2004, flannigan2000}.
\end{proposition}

\begin{novelbox}{Novel Contribution 3 --- Cumulative Vegetation Stress Index (CVSI)}
\begin{itemize}[leftmargin=1.5em,itemsep=2pt]
  \item No prior Indian national-scale forest fire study has introduced
        a temporally accumulating pre-fire NDVI stress metric.
  \item \citet{biswas2025} use NDVI at the analysis date only ---
        they have no memory of preceding months.
  \item CVSI quantifies multi-month fuel moisture deficit before
        ignition: a physically grounded causal feature.
  \item Optimal lag $k^{*}$ is selected via mutual information
        (Eq.~\ref{eq:mi_lag}) --- a rigorous information-theoretic
        criterion absent from all prior Indian fire studies.
  \item \textbf{References:} \citet{chuvieco2004}; \citet{flannigan2000}.
\end{itemize}
\end{novelbox}

% ==========================================================
\section{Feature Engineering Summary and Integration}
\label{sec:summary}
% ==========================================================

Table~\ref{tab:features_math} presents all seven NDVI-derived features
with their complete mathematical definitions.

\begin{table}[H]
\centering
\caption{Complete mathematical summary of NDVI-derived features}
\label{tab:features_math}
\small
\begin{tabularx}{\textwidth}{@{}clXc@{}}
\toprule
\textbf{ID} & \textbf{Feature} & \textbf{Defining Equation} &
\textbf{Novel?} \\
\midrule
F1 & QA NDVI &
  $\mathrm{NDVI}_{ij}^{(t)} = \mathrm{raw}_{ij}^{(t)} \times 10^{-4}$
  if $q_{ij}^{(t)} \in \{0,1\}$; else NaN & Partial \\[4pt]
F2 & Climatological mean &
  $\bar{\mu}_{ij}^{(m)} = \frac{1}{|\mathcal{Y}_m|}
  \sum_{y \in \mathcal{Y}_m} \mathrm{NDVI}_{ij}^{(y,m)}$ & No \\[4pt]
\rowcolor{lightgreen}
F3 & Anomaly $\delta$ &
  $\delta_{ij}^{(y,m)} = \mathrm{NDVI}_{ij}^{(y,m)} -
  \bar{\mu}_{ij}^{(m)}$ & \textbf{Yes} \\[4pt]
\rowcolor{lightgreen}
F4 & STL Trend $T(t)$ &
  $\mathrm{NDVI}_{ij}^{(t)} = T_{ij}^{(t)} + S_{ij}^{(t)} +
  R_{ij}^{(t)}$; LOESS smoother with period $p=12$ & \textbf{Yes} \\[4pt]
\rowcolor{lightgreen}
F5 & STL Residual $R(t)$ &
  $R_{ij}^{(t)} = \mathrm{NDVI}_{ij}^{(t)} - T_{ij}^{(t)} -
  S_{ij}^{(t)}$ & \textbf{Yes} \\[4pt]
\rowcolor{lightgreen}
F6 & M-K $\tau$ &
  $\tau = (P - Q) \,/\, \frac{1}{2}n(n-1)$; sig.\ if $p < 0.05$
  & \textbf{Yes} \\[4pt]
\rowcolor{lightgreen}
F7 & CVSI$(k^{*})$ &
  $\mathrm{CVSI}_{ij}(\tau_f, k^{*}) =
  \sum_{t=\tau_f - k^{*}}^{\tau_f - 1} \max(-\delta_{ij}^{(t)}, 0)$
  & \textbf{Yes} \\
\bottomrule
\end{tabularx}
\end{table}

These seven features replace the single time-averaged NDVI of
\citet{biswas2025} and expand the predictor space from 15 to 21
dimensions in the ML pipeline (RF, XGBoost, SVM + MaxEnt).
SHAP (SHapley Additive exPlanations) analysis \citep{lundberg2017}
is applied post-training to decompose each feature's marginal
contribution.

\subsection{Concrete Validation: Uttarakhand Fire Pixel}

Table~\ref{tab:example} compares what \citet{biswas2025} see versus
what this study extracts for a single fire pixel in Pauri Garhwal
district, Uttarakhand ($30.2^{\circ}$N, $78.8^{\circ}$E) ---
a confirmed high-fire zone from MODIS FIRMS data (April 2016).

\begin{table}[H]
\centering
\caption{Information extracted per methodology ---
         Uttarakhand fire pixel (April 2016)}
\label{tab:example}
\small
\begin{tabular}{@{}lll@{}}
\toprule
\textbf{Information} &
\textbf{Biswas et al.\ (2025)} &
\textbf{This Study} \\
\midrule
NDVI input &
  0.45 (single 20-year average) &
  309 monthly values decomposed \\
Long-term trend &
  Not detectable &
  $T$: $-0.003$ NDVI/year; net loss $= 0.075$ \\
Trend significance &
  Cannot test &
  $\tau = -0.31$, $p = 0.002$ --- significant browning \\
Pre-fire anomalies &
  Not detectable &
  $\delta_{\text{Jan}}=-0.03$, $\delta_{\text{Feb}}=-0.05$,
  $\delta_{\text{Mar}}=-0.04$ \\
CVSI ($k=3$) &
  Not available &
  $0.03 + 0.05 + 0.04 = 0.12$ (high stress) \\
\midrule
Model verdict &
  \textit{Moderate risk} &
  \textbf{HIGH FIRE RISK} \\
\bottomrule
\end{tabular}
\end{table}

% ==========================================================
\section{Conclusion}
\label{sec:conclusion}
% ==========================================================

This technical note has presented five novel extensions to the NDVI
preprocessing methodology of \citet{biswas2025}, summarised as follows:

\begin{enumerate}[leftmargin=2em, label=\textbf{NC\arabic*.}]
  \item \textbf{STL Decomposition} (Eqs.~\ref{eq:stl}--\ref{eq:residual}):
        pixel-wise separation of NDVI into trend $T$, seasonal $S$,
        and residual $R$ components over 309 monthly time steps.
  \item \textbf{Mann--Kendall Test} (Eqs.~\ref{eq:mk_tau}--\ref{eq:mk_decision}):
        non-parametric statistical validation of browning/greening
        trends at the $\alpha = 0.05$ level.
  \item \textbf{Monthly Anomaly} (Eq.~\ref{eq:anomaly}):
        removal of the seasonal cycle to expose inter-annual
        vegetation stress independent of time of year.
  \item \textbf{CVSI} (Eq.~\ref{eq:cvsi}):
        a novel multi-month pre-fire fuel moisture deficit index
        with information-theoretically optimal lag selection.
  \item \textbf{QA Filtering} (Eq.~\ref{eq:qa_filter}):
        pixel-level cloud and snow screening absent in \citet{biswas2025}.
\end{enumerate}

Together these transform 1 raw NDVI feature into 7 physically
interpretable, temporally dynamic predictors (F1--F7), expanding the
ML feature space from 15 to 21 dimensions relative to the baseline
study. The methodology is applicable at the national scale for India's
full claimed territory using 309 monthly MOD13A3 v061 files
(March 2000 -- December 2025).

% ==========================================================
% BIBLIOGRAPHY
% ==========================================================
\clearpage
\bibliographystyle{abbrvnat}
\bibliography{ndvi_novel_refs}

% ── Inline bibliography (if .bib file unavailable) ─────────
% Uncomment below and remove \bibliography{} above if needed:

% \begin{thebibliography}{99}
%
% \bibitem[Biswas et al.(2025)]{biswas2025}
% Biswas, U., Mahato, S., Joshi, P.K.\ (2025).
% Forest fire occurrence probability modelling using MaxEnt in India.
% \textit{Environmental Science and Pollution Research}, 32:4856--4878.
%
% \bibitem[Cleveland et al.(1990)]{cleveland1990stl}
% Cleveland, R.B., Cleveland, W.S., McRae, J.E., Terpenning, I.\ (1990).
% STL: A seasonal-trend decomposition procedure based on LOESS.
% \textit{Journal of Official Statistics}, 6(1):3--73.
%
% \bibitem[Mann(1945)]{mann1945}
% Mann, H.B.\ (1945).
% Nonparametric tests against trend.
% \textit{Econometrica}, 13(3):245--259.
%
% \bibitem[Kendall(1975)]{kendall1975}
% Kendall, M.G.\ (1975).
% \textit{Rank Correlation Methods}, 4th ed.
% Griffin, London.
%
% \bibitem[Chuvieco et al.(2004)]{chuvieco2004}
% Chuvieco, E., Cocero, D., Riano, D., Martin, P., Aguado, I.,
% de la Riva, J., Burgan, R.\ (2004).
% Combining NDVI and surface temperature for the estimation of live
% fuel moisture content in forest fire danger rating.
% \textit{Remote Sensing of Environment}, 92:322--331.
%
% \bibitem[Flannigan et al.(2000)]{flannigan2000}
% Flannigan, M.D., Stocks, B.J., Wotton, B.M.\ (2000).
% Climate change and forest fires.
% \textit{Science of the Total Environment}, 262:221--229.
%
% \bibitem[Lundberg and Lee(2017)]{lundberg2017}
% Lundberg, S.M., Lee, S.I.\ (2017).
% A unified approach to interpreting model predictions.
% \textit{Advances in Neural Information Processing Systems}, 30.
%
% \bibitem[Didan(2021)]{didan2021}
% Didan, K.\ (2021).
% MOD13A3 MODIS/Terra Vegetation Indices Monthly L3 Global
% 1km SIN Grid V061. NASA EOSDIS LP DAAC.
% \url{https://doi.org/10.5067/MODIS/MOD13A3.061}
%
% \end{thebibliography}

\end{document}
